%===============================================================================
\section{Objective~2: Application}
%-------------------------------------------------------------------------------
\begin{frame}{Bayesian model}
  \pbox{.75}{
  \begin{itemize}
    \item Implemented $f(z|s)$ equation in Stan
    \item Consider 5 candidate families for true $f(x)$:
       \\ \emph{Exponential, Gamma, Weibull, Scaled Beta, Scaled Kumaraswamy}
       \\ with log-uniform priors on parameters (\textcolor{cr}{$\aa,\bb,\pp,\tt$})
    \item Fit to 2 types of aggregate data:
    \begin{itemize}
      \item \textcolor{cd}{$\hat\E[z|s;N_i]$}: mean durations
      \item \textcolor{cd}{$\hat\P[\azb|s;N_i]$}: interval proportions (includes medians, quantiles, etc.)
    \end{itemize}
    where $N_i$ denotes the study sample size
    \item Keep 1000 samples after 500 burn-in iterations from 7 chains per family
    \item Final posterior: importance resample across all families
  \end{itemize}}
  \pbox{.25}{\centering\includegraphics[scale=1]{model}}
\end{frame}
%-------------------------------------------------------------------------------
\begin{frame}{Data sources}
  \pbox[t]{.5}{\paragraph{1.~Proof-of-concept}
    Aggregate data (RDS-adjusted) from\\
    313 female sex workers in Eswatini:\cite{Baral2014}
    \pause\par\bigskip\begin{tblr}[theme=slides]{colspec={rrc}}
{Duration\\(years)} & \SetCell[c=2]{c}{Population \%\\mean \quad(\CI)}\\
0--3     & 38.3 & (27.5,~49.1) \\
3--6     & 32.1 & (23.6,~40.7) \\
6--11    & 20.2 & (13.2,~27.1) \\
11+      &  9.4 &  (4.4,~14.4) \\
\end{tblr}
}
  \pbox[t]{.5}{\paragraph{2.~Regional synthesis}
     Aggregate data from\\
     58,249 female sex workers globally:
     \begin{itemize}
       \item 47 means
       \item 39 medians
       \item 126 quantiles
     \end{itemize}
     \pause
     representing 212 estimates from 82 studies\\
     reviewed by \citet{Fazito2012}}
     % 14 -- 15,277 respondents
\end{frame}
%-------------------------------------------------------------------------------
\begin{frame}{Proof-of-concept: fit vs data}
  \itergraphics[width=\tw]{demo.props.s}{4}{1,2,3,4}
\end{frame}
%-------------------------------------------------------------------------------
\begin{frame}{Proof-of-concept: parameter posteriors}
  \centering\includegraphics[width=.8\tw]{demo.pars.s}
\end{frame}
%-------------------------------------------------------------------------------
\begin{frame}{Regional synthesis: fit vs data}
  \centering\itergraphics[width=.9\tw]{meta.distr.s}{9}{1,2,3,4,5,6,7,8,9}
\end{frame}
%-------------------------------------------------------------------------------
\begin{frame}{Regional synthesis: infer shorter mean durations $\E[x]$ for all regions}
  \centering\includegraphics[width=.5\tw]{meta.forest.s}
  % Why: review methods ignored length-biased sampling \& undersampling of early sex work
\end{frame}
%-------------------------------------------------------------------------------
\begin{frame}{Shorter mean durations $\E[x]$: implications}
  \begin{itemize}
    \item \textbf{Turnover}: shorter durations means faster turnover
    \item \textbf{STI/HIV Incidence}: shorter durations means less time for infections to accumulate
       \\ \eg 10\% STI/HIV prevalence in sex work:
    \begin{itemize}
      \item with \z1-year mean duration \rarr high incidence
      \item with  10-year mean duration \rarr lower incidence
    \end{itemize}
    % \item \textbf{Onward transmission}: shorter durations mean less time for onward transmission via sex work
    \item \textbf{Lifetime exposure}: faster turnover means more women ever sold sex?
       \\ (well-known under-reporting \cite{Behanzin2013})
    \item \textbf{Program coverage}: faster turnover means constant influx of new sex workers to reach
       \\ especially for vulnerable early period\cite{Cheuk2020}
  \end{itemize}
\end{frame}
%-------------------------------------------------------------------------------
\begin{frame}{Regional synthesis: heavy-tail duration distributions}
  \hfill\includegraphics[width=.67\tw]{meta.distr.x}\hfill\null
  \par\rarr Infer 50+\% women sell sex for <~1~year in most regions
\end{frame}
%-------------------------------------------------------------------------------
\begin{frame}{Heavy-tail duration distributions: implications}
  \begin{itemize}
    \item \textbf{Exit rates}: decline with longer duration \rarr heavy-tail distribution
    \item \textbf{Population Heterogeneity}: could explain this
    \begin{itemize}
      \item within-person heterogeneity e.g. changes in self-identification / social context
      \item between-person heterogeneity e.g. selection effects (frailties)
    \end{itemize}
    \item \textbf{Modelling} this heterogeneity, need:
    \begin{itemize}
      \item early vs late sex work periods, and/or
      \item short vs long duration trajectories
    \end{itemize}
    \item \textbf{Programs}: critical to reach early sex work period
  \end{itemize}
\end{frame}
